\clearpage
\SourceLeaf{144}
\begin{center}
Monatliche Correspondenz zur Beförderung der Erd- und Himmelskunde. Herausgegeben von F.\ \textsc{von Zach}.\\
Erster Band. S.\,193. Gotha, 1800.
\end{center}

\begin{center}
\textls{Verbesserungen zum IV. Bande der Allg.\ Geogr.\ Ephemer.}
\end{center}

\begin{center}
. . . . . . . . . .
\end{center}

Ebendaselbst S.\,378 Nr.\,3 Zeile 9 statt $\frac{1}{150}$ Abplattung muss es heissen $\frac{1}{50}$ Abplattung. Zeile 12 statt der Worte »\emph{Der Unterschied von $\frac{1}{150}$ und $\frac{1}{187}$ ist in diesem Falle eben nicht erheblich}« kann man zu mehrerer Verständlichkeit folgendes setzen »\emph{Der Unterschied von $\frac{1}{150}$, [welche Abplattung] die Französischen Grad-Messer (A.\,G.\,E. IV.\,B.\ S.\,XXXVII der Einleitung und S.\,42) und $\frac{1}{187}$, die ich gefunden habe, ist in diesem Falle eben nicht erheblich}«.

\bigskip
\begin{center}
[\textsc{Schumacher} \emph{an} \textsc{Gauss}. \emph{Altona, 30. November 1831.}]
\end{center}

\begin{center}
. . . . . . . . . . .
\end{center}

\{Ich glaube Ihnen schon einmal gesagt zu haben, dass \textsc{Zach} in den Geographischen Ephemeriden (1799, October.\ p.\,378) einen Brief von Ihnen hat abdrucken lassen, in dem Sie offenbar die Methode der kleinsten Quadrate erwähnen, die Sie also damals schon \textsc{Zach} mitgetheilt haben. Sie sprechen von der französischen Gradmessung:

\begin{quote}
»Ich entdeckte diesen Fehler, indem ich meine Methode, von der ich \textls{Ihnen eine Probe gegeben habe}, anwandte« u.\,s.\,w.
\end{quote}

\textsc{Zach} bemerkt dabei: »Hiervon ein andermal«, das andere Mal ist aber nie gekommen. Da Sie die Resultate Ihrer Rechnung geben, so scheint es mir, ist es leicht zu zeigen, dass diese durch die Methode der kleinsten Quadrate abgeleitet sind. \textsc{Zach} lebt zudem noch, und hat gewiss Ihren Brief aufgehoben. Finden Sie es nicht der Mühe werth, endlich die Sache einmal, selbst gegen die mir vor allen widerlichen \textls{höflichen} Zweifel der Franzosen, unwidersprechlich abzumachen?\}

\begin{center}
. . . . . . . . . . . . .
\end{center}
